% GENERATED FILE. Edit canonical lessons, then run npm run generate:books.
% canonical-source-hash: fnv1a64:4838490d80e0d1d3
% canonical-lessons: PA-C244-pona, PA-C244-punna, PA-C244-nitarna, PA-C244-parosna, PA-C244-ghotna

\chapter{From Pot to Plate}
\label{ch:pa-a2-from-pot-to-plate}
\hlchaptermodality{244}

\begin{chapteropening}
\textbf{By the end of this chapter:} \emph{I can talk about the last steps of making a meal.}

Complete the last lesson of chapter 244: I can talk about the last steps of making a meal.
\end{chapteropening}

\section[\texorpdfstring{\pa{ਪੋਣਾ} (poṇā)}{ਪੋਣਾ (poṇā)}]{\pa{ਪੋਣਾ} (poṇā) — a straining cloth}

\label{lesson:PA-C244-pona}



\begin{quote}
Before the new one: say the Punjabi for to sift, then the Punjabi for to taste.
\end{quote}

\subsection*{You'll want to know: \pa{ਪੋਣਾ}}
\textbf{\pa{ਪੋਣਾ}} — \emph{poṇā} — \textquotedblleft{}a straining cloth\textquotedblright{}.

\begin{quote}
\pa{ਪੋਣਾ} \pa{ਚਿੱਟਾ} \pa{ਅਤੇ} \pa{ਸਾਫ਼} \pa{ਹੈ।}
\end{quote}

\begin{quote}
\emph{poṇā chiṭṭā ate sāf hai.}
\end{quote}

\begin{quote}
The straining cloth is white and clean.
\end{quote}

A masculine noun.

\begin{cousinweb}
\textbf{\pa{ਪੋਣਾ}} (\emph{poṇā}) is inherited from Sanskrit \emph{pavana}, \textquotedblleft{}sieve, strainer\textquotedblright{}; CDIAL lists Punjabi \emph{poṇā}, \textquotedblleft{}strainer, sieve\textquotedblright{} (CDIAL 7977).
\end{cousinweb}

\begin{grammarlens}[title={What you've built}]
One more word for this chapter.
\end{grammarlens}

\subsection*{Your turn}
\begin{itemize}
\raggedright
  \item \emph{Say it:} \emph{poṇā}
  \item \emph{Say it:} \emph{poṇā}, once more
  \item \emph{Say it:} \emph{poṇā chiṭṭā ate sāf hai.} — The straining cloth is white and clean.
  \item \emph{Recall:} say the Punjabi for to sift, then the Punjabi for to taste, then say \emph{poṇā} again
\end{itemize}

\begin{tcolorbox}[breakable,colback=teal!4,colframe=teal!35!black,title={Before you move on}]
What does \pa{ਪੋਣਾ} mean? (\textquotedblleft{}A straining cloth\textquotedblright{}.) Say it once more.
\end{tcolorbox}
\section[\texorpdfstring{\pa{ਪੁਣਨਾ} (puṇnā)}{ਪੁਣਨਾ (puṇnā)}]{\pa{ਪੁਣਨਾ} (puṇnā) — to strain}

\label{lesson:PA-C244-punna}



\begin{quote}
Before the new one: say the Punjabi for to taste, then the Punjabi for a straining cloth.
\end{quote}

\subsection*{You'll want to know: \pa{ਪੁਣਨਾ}}
\textbf{\pa{ਪੁਣਨਾ}} — \emph{puṇnā} — \textquotedblleft{}to strain\textquotedblright{}.

\begin{quote}
\pa{ਮੈਂ} \pa{ਪੋਣੇ} \pa{ਨਾਲ} \pa{ਚਾਹ} \pa{ਪੁਣਦੀ} \pa{ਹਾਂ।}
\end{quote}

\begin{quote}
\emph{maiṁ poṇe nāl chāh puṇdī hāṁ.}
\end{quote}

\begin{quote}
I strain the tea through a cloth.
\end{quote}

Past (I …): \pa{ਮੈਂ} \pa{ਚਾਹ} \pa{ਪੁਣੀ}, \emph{maiṁ chāh puṇī}.

\begin{cousinweb}
\textbf{\pa{ਪੁਣਨਾ}} (\emph{puṇnā}) is inherited from Sanskrit \emph{punāti}, \textquotedblleft{}purifies\textquotedblright{}; CDIAL lists Punjabi \emph{puṇnā}, \textquotedblleft{}to strain, darn\textquotedblright{} (CDIAL 8277).
\end{cousinweb}

\begin{grammarlens}[title={What you've built}]
One more word for this chapter.
\end{grammarlens}

\subsection*{Your turn}
\begin{itemize}
\raggedright
  \item \emph{Say it:} \emph{puṇnā}
  \item \emph{Say it:} \emph{puṇnā}, once more
  \item \emph{Say it:} \emph{maiṁ poṇe nāl chāh puṇdī hāṁ.} — I strain the tea through a cloth.
  \item \emph{Recall:} say the Punjabi for to taste, then the Punjabi for a straining cloth, then say \emph{puṇnā} again
  \item \emph{Write it:} copy \textbf{\pa{ਪੁਣਨਾ}} once, checking the whole word
\end{itemize}

\begin{tcolorbox}[breakable,colback=teal!4,colframe=teal!35!black,title={Before you move on}]
What does \pa{ਪੁਣਨਾ} mean? (\textquotedblleft{}To strain\textquotedblright{}.) Say it once more.
\end{tcolorbox}
\section[\texorpdfstring{\pa{ਨਿਤਾਰਨਾ} (nitārnā)}{ਨਿਤਾਰਨਾ (nitārnā)}]{\pa{ਨਿਤਾਰਨਾ} (nitārnā) — to pour off (water)}

\label{lesson:PA-C244-nitarna}



\begin{quote}
Before the new one: say the Punjabi for a straining cloth, then the Punjabi for to strain.
\end{quote}

\subsection*{You'll want to know: \pa{ਨਿਤਾਰਨਾ}}
\textbf{\pa{ਨਿਤਾਰਨਾ}} — \emph{nitārnā} — \textquotedblleft{}to pour off (water)\textquotedblright{}.

\begin{quote}
\pa{ਮੈਂ} \pa{ਚੌਲਾਂ} \pa{ਦਾ} \pa{ਪਾਣੀ} \pa{ਨਿਤਾਰਦੀ} \pa{ਹਾਂ।}
\end{quote}

\begin{quote}
\emph{maiṁ chaulāṁ dā pāṇī nitārdī hāṁ.}
\end{quote}

\begin{quote}
I pour off the water from the rice.
\end{quote}

Past (I …): \pa{ਮੈਂ} \pa{ਪਾਣੀ} \pa{ਨਿਤਾਰਿਆ}, \emph{maiṁ pāṇī nitāriā}.

\begin{cousinweb}
\textbf{\pa{ਨਿਤਾਰਨਾ}} (\emph{nitārnā}) is inherited from Sanskrit \emph{nistārayati}, \textquotedblleft{}delivers from\textquotedblright{}; CDIAL lists Punjabi \emph{nitārṇā}, \textquotedblleft{}to clarify\textquotedblright{} (CDIAL 7532). The sense \textquotedblleft{}pour off\textquotedblright{} is CDIAL's Lahnda gloss: \emph{nitāraṇ}, \textquotedblleft{}to skim, decant, pour off\textquotedblright{}.
\end{cousinweb}

\begin{grammarlens}[title={What you've built}]
One more word for this chapter.
\end{grammarlens}

\subsection*{Your turn}
\begin{itemize}
\raggedright
  \item \emph{Say it:} \emph{nitārnā}
  \item \emph{Say it:} \emph{nitārnā}, once more
  \item \emph{Say it:} \emph{maiṁ chaulāṁ dā pāṇī nitārdī hāṁ.} — I pour off the water from the rice.
  \item \emph{Recall:} say the Punjabi for a straining cloth, then the Punjabi for to strain, then say \emph{nitārnā} again
  \item \emph{Write it:} copy \textbf{\pa{ਨਿਤਾਰਨਾ}} once, checking the whole word
\end{itemize}

\begin{tcolorbox}[breakable,colback=teal!4,colframe=teal!35!black,title={Before you move on}]
What does \pa{ਨਿਤਾਰਨਾ} mean? (\textquotedblleft{}To pour off (water)\textquotedblright{}.) Say it once more.
\end{tcolorbox}
\section[\texorpdfstring{\pa{ਪਰੋਸਣਾ} (parosṇā)}{ਪਰੋਸਣਾ (parosṇā)}]{\pa{ਪਰੋਸਣਾ} (parosṇā) — to serve food}

\label{lesson:PA-C244-parosna}



\begin{quote}
Before the new one: say the Punjabi for to strain, then the Punjabi for to pour off.
\end{quote}

\subsection*{You'll want to know: \pa{ਪਰੋਸਣਾ}}
\textbf{\pa{ਪਰੋਸਣਾ}} — \emph{parosṇā} — \textquotedblleft{}to serve food\textquotedblright{}.

\begin{quote}
\pa{ਮਾਤਾ} \pa{ਮਹਿਮਾਨਾਂ} \pa{ਨੂੰ} \pa{ਖਾਣਾ} \pa{ਪਰੋਸਦੀ} \pa{ਹੈ।}
\end{quote}

\begin{quote}
\emph{mātā mahimānāṁ nū̃ khāṇā parosdī hai.}
\end{quote}

\begin{quote}
Mother serves food to the guests.
\end{quote}

Past (I …): \pa{ਮੈਂ} \pa{ਖਾਣਾ} \pa{ਪਰੋਸਿਆ}, \emph{maiṁ khāṇā parosiā}.

\begin{cousinweb}
\textbf{\pa{ਪਰੋਸਣਾ}} (\emph{parosṇā}) goes back to Sanskrit \emph{parivēṣati}, \textquotedblleft{}attends, serves food\textquotedblright{}, as a word taken from Hindi: CDIAL lists Punjabi \emph{parosṇā}, \textquotedblleft{}to offer food to a guest\textquotedblright{}, as borrowed from Hindi \emph{parosnā}; the Punjabi form inherited directly is \emph{parīhṇā}, \textquotedblleft{}to serve a meal\textquotedblright{} (CDIAL 7888).
\end{cousinweb}

\begin{grammarlens}[title={What you've built}]
One more word for this chapter.
\end{grammarlens}

\subsection*{Your turn}
\begin{itemize}
\raggedright
  \item \emph{Say it:} \emph{parosṇā}
  \item \emph{Say it:} \emph{parosṇā}, once more
  \item \emph{Say it:} \emph{mātā mahimānāṁ nū̃ khāṇā parosdī hai.} — Mother serves food to the guests.
  \item \emph{Recall:} say the Punjabi for to strain, then the Punjabi for to pour off, then say \emph{parosṇā} again
  \item \emph{Write it:} copy \textbf{\pa{ਪਰੋਸਣਾ}} once, checking the whole word
\end{itemize}

\begin{tcolorbox}[breakable,colback=teal!4,colframe=teal!35!black,title={Before you move on}]
What does \pa{ਪਰੋਸਣਾ} mean? (\textquotedblleft{}To serve food\textquotedblright{}.) Say it once more.
\end{tcolorbox}
\section[\texorpdfstring{\pa{ਘੋਟਣਾ} (ghoṭṇā)}{ਘੋਟਣਾ (ghoṭṇā)}]{\pa{ਘੋਟਣਾ} (ghoṭṇā) — to mash; to grind smooth}

\label{lesson:PA-C244-ghotna}



\begin{quote}
Before the new one: say the Punjabi for to pour off, then the Punjabi for to serve food.
\end{quote}

\subsection*{You'll want to know: \pa{ਘੋਟਣਾ}}
\textbf{\pa{ਘੋਟਣਾ}} — \emph{ghoṭṇā} — \textquotedblleft{}to mash; to grind smooth\textquotedblright{}.

\begin{quote}
\pa{ਦਾਦੀ} \pa{ਸਾਗ} \pa{ਬਹੁਤ} \pa{ਦੇਰ} \pa{ਘੋਟਦੀ} \pa{ਹੈ।}
\end{quote}

\begin{quote}
\emph{dādī sāg bahut der ghoṭdī hai.}
\end{quote}

\begin{quote}
Grandmother mashes the greens for a long time.
\end{quote}

Past (I …): \pa{ਮੈਂ} \pa{ਸਾਗ} \pa{ਘੋਟਿਆ}, \emph{maiṁ sāg ghoṭiā}.

\begin{cousinweb}
\textbf{\pa{ਘੋਟਣਾ}} (\emph{ghoṭṇā}) is inherited from Sanskrit \emph{ghaṭṭayati}, \textquotedblleft{}rubs, touches, shakes\textquotedblright{}; CDIAL lists Punjabi \emph{ghoṭṇā}, \textquotedblleft{}to grind, rub, smoothe\textquotedblright{} (CDIAL 4417).
\end{cousinweb}

\begin{grammarlens}[title={What you've built}]
One more word for this chapter.
\end{grammarlens}

\subsection*{Your turn}
\begin{itemize}
\raggedright
  \item \emph{Say it:} \emph{ghoṭṇā}
  \item \emph{Say it:} \emph{ghoṭṇā}, once more
  \item \emph{Say it:} \emph{dādī sāg bahut der ghoṭdī hai.} — Grandmother mashes the greens for a long time.
  \item \emph{Recall:} say the Punjabi for to pour off, then the Punjabi for to serve food, then say \emph{ghoṭṇā} again
  \item \emph{Write it:} copy \textbf{\pa{ਘੋਟਣਾ}} once, checking the whole word
\end{itemize}

\begin{tcolorbox}[breakable,colback=teal!4,colframe=teal!35!black,title={Before you move on}]
What does \pa{ਘੋਟਣਾ} mean? (\textquotedblleft{}To mash; to grind smooth\textquotedblright{}.) Say it once more.
\end{tcolorbox}
